%% sections/renzhi.tex: Renzhi's slides (comparison across methods; the LLM designer loop). OWNER: Renzhi. Only Renzhi edits this file.
%% - Figures go in figures/renzhi/ and are included as \includegraphics[width=...]{renzhi/<file>}.
%% - Need a macro of your own? Define it here with a renzhi prefix (\newcommand{\renzhiFoo}{...}), never in preamble.tex.
%% - Build just your part: latexmk -pdf solo_renzhi.tex   (or compile solo_renzhi.tex in Overleaf).
%% - Backup slides: put them at the end of this file as \begin{frame}[noframenumbering]{Backup: ...}.
%% - The section's title and cover slide are set in main.tex; ask Renzhi to rename it.


%% Numbers: rl/PLAN.md (references v4, test references, the ladder, live validation, robustness R1-R3), the headroom
%% check (rl/PLAN.md on branch rl/es-fitness, 2026-11-25) and doc/final/final.tex. SKY130 unless marked GF180.

\begin{frame}{Mapping gets hard fast: the knob ladder}
  \begin{columns}[c]
    \begin{column}{0.55\textwidth}
      \centering
      \begin{tikzpicture}[x=1mm, y=38mm]
        \foreach \y in {0, 0.5, 1} {
          \draw[uscgray!30] (0,\y) -- (62,\y);
          \node[left, font=\scriptsize, text=uscgray] at (0,\y) {\y};
        }
        \foreach \x/\v/\e/\k in {4/0.90/0.052/8, 19/0.60/0.113/15, 34/0.32/0.121/22, 49/0.15/0.176/29} {
          \fill[usccardinal] (\x,0) rectangle ++(9,\v);
          \node[above, font=\scriptsize] at (\x+4.5,\v) {\v};
          \node[below, font=\scriptsize] at (\x+4.5,0) {\k};
        }
        \draw[uscgray] (0,0) -- (62,0);
        \node[font=\scriptsize, text=uscgray] at (31,-0.16) {knobs};
        \node[rotate=90, font=\scriptsize, text=uscgray] at (-9,0.5) {\passfail{}};
      \end{tikzpicture}
    \end{column}
    \begin{column}{0.42\textwidth}
      \small
      \begin{itemize}
        \item One amplifier, the same specs, 8 to 29 free knobs; the \rulename{}, 600 simulations.
        \item \passfail{} falls from 0.90 to 0.15.
        \item \prederr{} rises from 0.052 to 0.176.
        \item The room for a better explorer is in the larger circuits.
      \end{itemize}
    \end{column}
  \end{columns}
\end{frame}

\begin{frame}{The bar on the test sets}
  \centering\small
  The \rulename{}, 10 seeds; \prederr{} at 50 / 200 / 600 simulations, \passfail{} at 600

  \vspace{2mm}
  \begin{tabular}{@{}lcc@{}}
    \toprule
    Test set & \prederr{} & \passfail{} \\
    \midrule
    Familiar size, 7--8 knobs & 0.607 / 0.306 / 0.194 & 0.494 \\
    Larger, 11--20 knobs & 0.429 / 0.309 / 0.237 & 0.251 \\
    Ladder, 8 / 22 / 29 knobs & 0.347 / 0.172 / 0.116 & 0.454 \\
    \midrule
    GF\,180, 3 circuits (23 tasks) & 0.597 / 0.247 / 0.139 & 0.672 \\
    \bottomrule
  \end{tabular}

  \vspace{3mm}
  \begin{minipage}{0.85\linewidth}\footnotesize
    \begin{itemize}
      \item Greedy is within 0.01 of the rule everywhere: the two fixed rules tie on stored pools.
      \item The map's chance that any design passes separates achievable from impossible specs perfectly.
    \end{itemize}
  \end{minipage}
\end{frame}

\begin{frame}{Headline: every method against the \rulename{} (test)}
  \centering\small
  Difference from the rule, 95\% interval over tasks; \prederr{}: lower is better, \passfail{}: higher is better

  \vspace{2mm}
  \begin{tabular}{@{}lcc@{}}
    \toprule
    Method & \prederr{} & \passfail{} \\
    \midrule
    \multicolumn{3}{@{}l}{\textit{Learned}} \\
    \quad ES (Taman) & \todo{test} & \todo{test} \\
    \quad PPO; PPO, discount 0; Milestone 1 PPO (Samir) & \todo{test} & \todo{test} \\
    \quad DQN (Hungyu) & \todo{test} & \todo{test} \\
    \multicolumn{3}{@{}l}{\textit{Baselines}} \\
    \quad GP straddle (Samir) & \todo{test} & \todo{test} \\
    \quad NSGA-II, MAP-Elites (Hungyu) & \todo{test} & \todo{test} \\
    \quad Greedy & within $\pm 0.01$ & within $\pm 0.01$ \\
    \bottomrule
  \end{tabular}

  \vspace{2mm}
  \footnotesize Test runs Nov 30 -- Dec 3 on the models frozen Nov 29.
\end{frame}

\begin{frame}{Stored pools against live simulation}
  \begin{columns}[c]
    \begin{column}{0.5\textwidth}
      \centering\small
      The \rulename{} on the 15-knob validation circuit, 600 simulations

      \vspace{2mm}
      \begin{tabular}{@{}lcc@{}}
        \toprule
        & stored & live \\
        \midrule
        \passfail{} at 600 & 0.732 & 0.818 \\
        passing designs found & 1$\times$ & 3.2$\times$ \\
        greedy $-$ rule, \passfail{} & $+0.003$ & $-0.031$ \\
        \bottomrule
      \end{tabular}
    \end{column}
    \begin{column}{0.47\textwidth}
      \small
      \begin{itemize}
        \item Live, every pick is a fresh design, simulated: 28,800 simulations.
        \item Stored candidates run out late in a run: the rule reaches 92\% of the pool's passing candidates by 600.
        \item Live, greedy falls behind the rule, where the stored run shows a tie.
        \item The map improves only moderately: the number of simulations, not the pool, mostly limits it.
      \end{itemize}
    \end{column}
  \end{columns}
\end{frame}

\begin{frame}{How much better could any picking rule do?}
  \begin{columns}[T]
    \begin{column}{0.5\textwidth}
      \small
      A \alert{one-step lookahead oracle} reads the true result of every offered candidate and picks the one that most
      improves the map, at every step. 19 validation tasks, 2 seeds.

      \vspace{2mm}
      \centering
      \begin{tabular}{@{}lcc@{}}
        \toprule
        Oracle $-$ rule & \passfail{} & \prederr{} \\
        \midrule
        improves margin & $-0.037$ & $\mathbf{-0.034}$ \\
        improves IoU & $-0.047$ & $+0.011$ \\
        \bottomrule
      \end{tabular}

      \vspace{1mm}
      \footnotesize Whole run; both \passfail{} intervals exclude 0.
    \end{column}
    \begin{column}{0.47\textwidth}
      \small
      \begin{itemize}
        \item Even with the truth, greedy improvement of the map \alert{loses to the rule} on \passfail{}.
          \item They find as many passing designs (recall 0.70--0.79, rule 0.75) but mark more wrong (precision
          0.78--0.83, rule 0.89).
        \item A per-pick reward does not match \passfail{}: the rule's uncertain picks pay off over many picks, which
          is why no reward or fitness change beat it.
      \end{itemize}
    \end{column}
  \end{columns}
\end{frame}

\begin{frame}{Transfer}
  \begin{columns}[T]
    \begin{column}{0.5\textwidth}
      \textbf{Spec to spec}
      \small
      \begin{itemize}
        \item A revised spec on the same circuit starts from the designs already simulated for the old one.
        \item The rule, 296 spec pairs: \prederr{} $-0.196$ [$-0.256$, $-0.150$], \passfail{} $+0.237$ [$+0.206$,
          $+0.265$] against starting fresh.
        \item Learned models: \todo{test}
      \end{itemize}
    \end{column}
    \begin{column}{0.47\textwidth}
      \textbf{SKY130 to GF\,180}
      \small
      \begin{itemize}
        \item 3 circuits, 23 tasks; the frozen SKY130 models run zero-shot.
        \item The same circuits work far more often on GF\,180 (3.3\,V supply): 15--61\% of their pools against
          4--49\%.
        \item Results: \todo{test}
      \end{itemize}
    \end{column}
  \end{columns}
\end{frame}

\begin{frame}{Why a region: robustness}
  \begin{columns}[c]
    \begin{column}{0.52\textwidth}
      \small
      One thin spec per validation circuit; each nominally passing design re-simulated at 5 MOSFET corners
      $\times$ 3 temperatures, resistor and capacitor corners, and 50 mismatch samples (15,387 simulations).
      \begin{itemize}
        \item Only \alert{0--16\%} of nominally passing designs pass at every corner and temperature.
        \item Mismatch alone keeps 75--99\%.
        \item Survival follows the \alert{nominal margin} (Spearman 0.59--0.73), not the distance to a failing design
          in knob space.
      \end{itemize}
    \end{column}
    \begin{column}{0.45\textwidth}
      \begin{block}{What it means}
        \small
        A first passing design is usually not a robust one. A map with predicted margins says which designs are
        worth keeping, which a single found design cannot.
      \end{block}
    \end{column}
  \end{columns}
\end{frame}

\begin{frame}{The map for the designer}
  \begin{columns}[T]
    \begin{column}{0.5\textwidth}
      \small
      Each simulated circuit gets a short \alert{map summary} in the designer's round report:
      \begin{itemize}
        \item the estimated feasible share, with a range;
        \item the chance that any design passes;
        \item the target that limits the region;
        \item the knob ranges that work, and the most interior design.
      \end{itemize}
    \end{column}
    \begin{column}{0.47\textwidth}
      \small
      \textbf{Accurate from 200 simulations} (13 validation tasks $\times$ 3 seeds):
      \begin{itemize}
        \item feasible share: median ratio to the truth about 1 (1.5--1.8$\times$ at 50);
        \item limiting target right on 82--92\% of tasks;
        \item achievable and impossible specs separated perfectly;
        \item working knob ranges overlap the true ones by about 0.75.
      \end{itemize}
    \end{column}
  \end{columns}
\end{frame}

%% ------------------------------------------------------------------------------------------------------------------
%% The LLM designer (SKY130 only; stand-in compiler). Numbers from the menu enumeration and the model runs of
%% 2026-11-24/25 (llm/PLAN.md log; sandbox/llm_enum/STATUS.md, sandbox/llm_longer/STATUS.md).
\tikzset{
  renzhiBox/.style={draw=usccardinal, fill=usccardinal!6, rounded corners=2pt, thick, align=center,
                    font=\footnotesize, text width=27mm, minimum height=12mm, inner sep=2pt},
  renzhiSim/.style={renzhiBox, draw=uscgray, fill=uscgray!8},
  renzhiArr/.style={-{Stealth[length=2mm]}, thick, uscgray},
}

\begin{frame}{The LLM designer: the loop}
  \centering
  \begin{tikzpicture}
    \node[renzhiBox] (spec) at (0,0)        {\textbf{Spec}\\ gain, bandwidth, load};
    \node[renzhiBox] (pick) at (36mm,0)     {\textbf{Designer (LLM)}\\ picks 3 menu entries};
    \node[renzhiSim] (comp) at (72mm,0)     {\textbf{Compile}\\ the circuit change};
    \node[renzhiSim] (scr)  at (108mm,0)    {\textbf{Quick screen}\\ random sizings};
    \node[renzhiSim] (srch) at (108mm,-24mm) {\textbf{Mapping search}\\ 200 designs, \rulename};
    \node[renzhiSim] (pm)   at (72mm,-24mm)  {\textbf{Phase margin}\\ $\geq 45^\circ$ if compensated};
    \node[renzhiBox] (rep)  at (36mm,-24mm)  {\textbf{Round report}\\ with the map summary};
    \node[renzhiBox] (call) at (0,-24mm)     {\textbf{Call}\\ accept / replan / reject};
    \draw[renzhiArr] (spec) -- (pick);
    \draw[renzhiArr] (pick) -- (comp);
    \draw[renzhiArr] (comp) -- (scr);
    \draw[renzhiArr] (scr) -- (srch);
    \draw[renzhiArr] (srch) -- (pm);
    \draw[renzhiArr] (pm) -- (rep);
    \draw[renzhiArr] (rep) -- (call);
    \draw[renzhiArr, usccardinal] (call.north east) -- node[right, font=\scriptsize, text=usccardinal, pos=0.4] {replan}
      (pick.south west);
  \end{tikzpicture}

  \vspace{4mm}
  \begin{minipage}{0.9\linewidth}\footnotesize
    \textbf{Menu:} every catalogue circuit the designer may see, as it is or one change away: a degeneration
    resistor, a cascode, or a second gain stage. Up to 3 rounds; it may reject a spec only after 3 rounds in which
    nothing met it.
  \end{minipage}
\end{frame}

\begin{frame}{The LLM designer: 4 specs beyond the catalogue}
  \centering\small\setlength{\tabcolsep}{4pt}
  \begin{tabular}{@{}lcccc@{}}
    \toprule
    Spec & Meet & Random & Sonnet (2 draws) & Qwen 14B \\
    \midrule
    \texttt{miller2 cl30.01}      & 1 of 28 & 0.33 & accept, 2,152 (both) & accept, 2,936 \\
    \texttt{miller2 cl1000.01}    & 1 of 28 & 0.32 & accept, 1,936 (both) & accept, 2,184 \\
    \texttt{miller2 cl1000.02}    & 0 of 28 & 0    & reject, 2,872/2,576 & reject, 3,064 \\
    \texttt{ota\_sym\_casc cl30.03} & 3 of 33 & 0.62 & accept, 1,672 (both) & accept, 2,624 \\
    \bottomrule
  \end{tabular}

  \vspace{3mm}
  \begin{minipage}{0.92\linewidth}\footnotesize
    \begin{itemize}
      \item No catalogue circuit, as it is, meets any of the four (14 circuits hidden from the designer).
      \item Meet: menu entries with a design that meets the spec. Random: a random designer's chance of meeting it
        within 3 rounds. Models: the call and the simulations used.
      \item The entry that meets: a five-transistor OTA with a second stage, with 45$^\circ$ phase margin or more.
      \item Qwen 7B met none of the four. SKY130, stand-in compiler.
    \end{itemize}
  \end{minipage}
\end{frame}

\begin{frame}{The LLM designer against a random designer}
  \centering
  \begin{tikzpicture}[x=1mm, y=32mm]
    \foreach \y in {0, 0.5, 1} {
      \draw[uscgray!30] (0,\y) -- (124,\y);
      \node[left, font=\scriptsize, text=uscgray] at (0,\y) {\y};
    }
    \node[rotate=90, font=\scriptsize, text=uscgray] at (-9,0.5) {chance of meeting in 3 rounds};
    \foreach \x/\r/\s/\lab in {6/0.33/1/{miller2\\cl30.01}, 37/0.32/1/{miller2\\cl1000.01},
                               68/0/0/{miller2\\cl1000.02}, 99/0.62/1/{ota\_sym\_casc\\cl30.03}} {
      \fill[uscgray!55] (\x,0) rectangle ++(8,\r);
      \fill[usccardinal] (\x+9,0) rectangle ++(8,\s);
      \node[above, font=\scriptsize] at (\x+4,\r) {\r};
      \node[below, font=\scriptsize, align=center] at (\x+8.5,0) {\lab};
    }
    \foreach \x in {6, 37, 99} { \node[above, font=\scriptsize, text=usccardinal] at (\x+13,1) {2 of 2}; }
    \node[font=\scriptsize, align=center, text=uscgray] at (76.5,0.3) {no entry\\meets it:\\not found};
    \draw[uscgray] (0,0) -- (124,0);
    \fill[uscgray!55] (2,1.2) rectangle ++(4,0.08);  \node[right, font=\scriptsize] at (7,1.24) {random designer};
    \fill[usccardinal] (40,1.2) rectangle ++(4,0.08); \node[right, font=\scriptsize] at (45,1.24) {Sonnet (met in 2 of 2 draws)};
  \end{tikzpicture}

  \vspace{2mm}
  \begin{minipage}{0.92\linewidth}\footnotesize
    Random designer: 3 untried menu entries a round, 10,000 seeded orders; it meets all three about 7\% of the time.
    When both meet, the simulations are about the same: Sonnet 1,672--2,152, random 1,884--2,190.
  \end{minipage}
\end{frame}

\begin{frame}{The LLM designer: takeaways}
  \begin{itemize}
    \item \textbf{Choosing the change is the value.} Sonnet read ``gain too low'' after round 1, added a second stage,
      and met all 3 achievable specs in both draws; a random designer meets all three about 7\% of the time.
    \item \textbf{Not a cheaper search.} When both meet, the simulations are about the same. Its reject of the
      fourth spec holds as ``not found'': 10$\times$ the search on the closest entry still met nothing.
    \item \textbf{A 14B open model is enough to drive the loop:} 3 of 4, every reply valid (7B: 0 of 4), with about a
      third more simulations than Sonnet.
    \item \textbf{Caveat:} 3 achievable specs, each met by one menu entry: a demonstration, not a general result.
  \end{itemize}
\end{frame}
